\section{Discussion}
\label{sec:discussion}

The results reported in Section~\ref{subsec:results} place reinforcement
learning significantly above the random policy but also significantly
below the hand-tuned threshold heuristic, with a reward gap of roughly
$13$ units and a Welch's $t$-test $p$-value below $10^{-4}$. This section steps back from the raw numbers and interrogates what
they allow us to claim---and what they do not---about the proposed approach,
then identifies the limitations of the current setup and the directions for
improvement.


\subsection{Positioning of the contribution}
\label{subsec:positioning}

Before discussing the empirical gap in detail, it is useful to restate
what this paper does and does not claim. We do not contribute a new RL
algorithm: every component of our pipeline (PPO, GAE, observation
normalization, advantage normalization, entropy bonus) is taken
off-the-shelf from prior work~\cite{schulman2017ppo,schulman2016gae,
andrychowicz2021matters}. The contribution is applicative and
methodological in a narrower sense: we isolate a business
decision-making problem whose RL formulation has not, to our knowledge,
been published; we build the smallest non-trivial simulator capable of
carrying that decision; and we subject the resulting pipeline to a
statistically sound comparison against two natural baselines---a random
policy and a rule-based heuristic whose only parameter has itself been
swept---with explicit confidence intervals and a significance test.
Under this scope, the empirical findings reported below should be read
as an \emph{upper bound on the expected performance of a minimal,
honestly reported, reproducible PPO pipeline on this problem class},
rather than as a claim about the best achievable performance in
production.


\subsection{Why PPO outperforms the random policy}
\label{subsec:why-ppo-beats-random}

The gap between PPO and the random policy does not merely reflect a
difference in immediate expected reward; it stems from the fact that the
agent learns to \emph{structure} its decisions around the state. Three
mechanisms drive this behavior.

First, the agent identifies a \emph{semantics} of coverage status: it learns
that renewing is of little value---or even detrimental---while coverage is
still largely valid, and that it becomes critical as expiration approaches.
This semantics is not provided upfront: it is extracted exclusively from the
reward signal, through observed trajectories.

Second, the agent learns a \emph{temporal credit assignment}: the GAE
estimator propagates the penalties of expiration and of uncovered failure
back to the decisions that made them possible several steps earlier. The
random policy, which does not adapt to the current state, cannot benefit
from this information and accumulates penalties mechanically.

Third, the agent \emph{conditions} its decisions on the full state, which
lets it implicitly weigh the cost of a renewal against the contract-cost
level and the risk level. The random policy treats all situations
interchangeably and is necessarily penalized by this lack of specificity.


\subsection{Why the heuristic policy can still win}
\label{subsec:why-heuristic-wins}

The slight numerical advantage of the heuristic over PPO does not reflect a
defect of the reinforcement-learning approach; it reflects the \emph{alignment}
between the structure of the considered problem and the form of the
heuristic rule.

The analysis in Section~\ref{subsec:results} already identified three primary
contributors to this gap: the fit between the heuristic's inductive bias and
the optimal policy, the function-approximation error of the neural network,
and the finite training horizon.

It is useful to add a more general reading of this phenomenon. In a
problem whose optimal policy is exactly representable as a deterministic,
one-dimensional rule, the RL agent enters an asymmetric game: it must
\emph{rediscover}, through exploration, a policy that the heuristic
possesses from step one. This asymmetry mechanically vanishes as the
structure of the problem becomes richer and the information required for
the decision exceeds what a simple rule can encode.

The sensitivity analysis in Section~\ref{subsec:sensitivity} sharpens
this reading, but also tempers it. The heuristic used in the main
comparison happens to sit on the plateau of its threshold sweep
(thresholds $3$ and $5$ produce identical rewards), so
Table~\ref{tab:results} reports the comparison against a near-optimal
instance of a one-parameter rule. Yet even the worst heuristic
threshold tested ($15$) remains above PPO: the rule family is robust
to threshold mis-specification, whereas PPO in its current form is
not. In a production setting this robustness of simple rules would be
a further argument in their favour; the extensions listed in
Section~\ref{subsec:improvements} are designed precisely to move the
evaluation into regimes where the rule's simplicity becomes a
liability rather than an asset.


\subsection{Limitations of the current setup}
\label{subsec:limitations}

Several limitations bound the scope of the reported conclusions.

\paragraph{Simulator dependence.}
Training and evaluation are conducted exclusively in a simulated environment.
Distributions over contract durations, cost levels and risk, as well as
failure probabilities, are defined by fixed parameters that have not been
matched against Philips France's real installed base. Any bias between these
distributions and reality will mechanically propagate to the learned policy.

\paragraph{Single-equipment framing.}
The MDP models a single piece of equipment at a time, without interactions
between contracts. This choice simplifies analysis and reproducibility but
elides budget couplings, volume constraints and portfolio effects that lie
at the heart of the original operational problem.

\paragraph{Sensitivity to reward shaping.}
The reward function involves several components whose coefficients are set
manually. The threshold heuristic closely matches the maximum of this
shaping. A revision of the relative weights may shift the optimum and
materially change the ranking of policies.

\paragraph{Absence of out-of-distribution evaluation.}
Policies are only evaluated on episodes drawn from the same distribution
as those used for training. No experiment quantifies the robustness of
PPO to a distribution shift, despite such shifts being inevitable in
production. The heuristic's robustness to a change in the renewal
threshold is characterised in Section~\ref{subsec:sensitivity}, but the
analogous sensitivity of PPO to environmental parameters (reward
weights, failure rates, cost levels) is not yet reported and is an
explicit item on the agenda of Section~\ref{subsec:improvements}.

\paragraph{Limited interpretability.}
The learned policy is represented by a neural network whose decisions are
not natively explainable. This opacity is a barrier to adoption by business
teams, who must justify coverage decisions to internal and external
stakeholders.

\paragraph{No hard operational constraints.}
Contractual budget, exclusivities, renewal eligibility windows and
contractual obligations do not appear explicitly in the MDP. In a production
deployment, such constraints must be \emph{guaranteed}, not merely
encouraged by the reward.


\subsection{Paths for improvement}
\label{subsec:improvements}

The limitations above suggest several concrete directions.

\paragraph{Extension to the full installed base.}
Moving from the single-equipment framing to the simultaneous management of a
heterogeneous portfolio exposes couplings (global budget, volume
constraints) and risk asymmetries that the single-threshold heuristic cannot
exploit. A per-equipment factorized policy with a shared critic---in the
spirit of the formulation sketched in Section~\ref{subsec:ppo}---constitutes
a natural first step.

\paragraph{Offline learning from historical data.}
Once available, the operational traces of Philips France may serve as the
basis for offline RL, avoiding any interaction with the production system.
Algorithms such as Conservative Q-Learning or Implicit Q-Learning allow
exploiting such data while controlling extrapolation risk.

\paragraph{Constrained learning.}
Budgetary and contractual constraints may be enforced explicitly, either
via action masking (prohibition at the softmax level of infeasible actions)
or through constrained-RL algorithms such as Constrained PPO and Lagrangian
methods.

\paragraph{Robustness and uncertainty quantification.}
Evaluation under distribution shift, training over an ensemble of
parametrically diverse simulators, and the use of Bayesian policies would
characterize and mitigate the model's sensitivity to gaps between
simulation and production.

\paragraph{Interpretability.}
Several levers may reduce policy opacity: distillation of the learned policy
into a decision tree or a set of rules, analysis of input-feature
importance, or hybrid formulations combining an interpretable rule with a
learned residual.

\paragraph{Comparison with stochastic optimization.}
Finally, confronting PPO with classical optimization methods---multi-stage
stochastic programming, rolling-horizon integer programs---would complement
the study with a comparison against a global policy obtained directly from
an explicit model of the problem.

\bigskip
\noindent
These directions converge toward a more demanding experimental framework in
which the simplifying assumptions of the present work are lifted one by one.
They form the core of the research program outlined in
Section~\ref{sec:conclusion}.
